\documentclass[11pt,a4paper]{article}

% Packages
\usepackage[utf8]{inputenc}
\usepackage[T1]{fontenc}
\usepackage{amsmath,amssymb}
\usepackage{graphicx}
\usepackage{booktabs}
\usepackage{hyperref}
\usepackage{cleveref}
\usepackage{natbib}
\usepackage{xcolor}
\usepackage{multirow}
\usepackage{adjustbox}
\usepackage{enumitem}
\usepackage[margin=1in]{geometry}
\usepackage{microtype}

% Metadata
\title{CultureDx: Culture-Adaptive Diagnostic Multi-Agent System\\
       with Evidence Grounding for Chinese Psychiatric\\
       Differential Diagnosis}

\author{Yu-Ning Tsao \\
  National Yunlin University of Science and Technology \\
  \texttt{M11237033@yuntech.edu.tw}}

\date{\today}

% Document
\begin{document}

\maketitle

\begin{abstract}
Large language models (LLMs) show promise for psychiatric differential diagnosis, but their parametric knowledge of non-Western clinical presentations remains weak. Chinese psychiatric interviews feature pervasive somatic idioms of distress (e.g., ``chest tightness'' for anxiety, ``head dizziness'' for depression) that map poorly to ICD-10 criteria without explicit cultural adaptation. We present CultureDx, a culture-adaptive multi-agent system (MAS) that decomposes clinical reasoning into specialized agent roles: triage routing, per-disorder criterion checking against ICD-10 criteria, deterministic logic evaluation, and statistical confidence calibration. Our evidence pipeline includes a 150-entry Chinese somatization ontology, temporal extraction, scope-aware negation detection, and hybrid dense-sparse retrieval. Experiments on LingxiDiag-16K and MDD-5k demonstrate that evidence-grounded MAS outperforms single-LLM baselines for Chinese psychiatric diagnosis, with particular gains in comorbidity detection and culture-specific symptom attribution.
% TODO: Add specific numbers once experiments complete
\end{abstract}

% Sections
\section{Introduction}
\label{sec:introduction}
% Introduction
% Motivation: LLMs lack Chinese psychiatric knowledge
% Contribution: Culture-adaptive MAS with evidence grounding

% TODO: Write introduction


\section{Related Work}
\label{sec:related-work}
% Related Work
% - LLMs for clinical NLP and psychiatric diagnosis
% - Multi-agent systems for medical reasoning
% - Chinese clinical NLP and somatization
% - Evidence-grounded diagnosis systems

% TODO: Write related work


\section{Method}
\label{sec:method}
% Method
% 3.1 System Architecture (HiED pipeline)
% 3.2 Evidence Pipeline (extraction, somatization, retrieval)
% 3.3 Diagnostic Logic (ICD-10 thresholds, calibration)
% 3.4 Comorbidity Resolution

% TODO: Write method section


\section{Experimental Setup}
\label{sec:experiments}
% Experimental Setup
% 4.1 Datasets (LingxiDiag-16K, MDD-5k)
% 4.2 Models (Qwen3-32B-AWQ, Qwen3-8B finetuned)
% 4.3 Baselines and Ablations
% 4.4 Evaluation Metrics

% TODO: Write experimental setup


\section{Results}
\label{sec:results}
% Results
% 5.1 Main Results (HiED vs baselines)
% 5.2 Ablation Study (evidence, somatization, reasoning)
% 5.3 Cross-Dataset Generalization
% 5.4 Comorbidity Detection

% TODO: Write results with tables from paper/tables/


\section{Analysis}
\label{sec:analysis}
% Analysis
% 6.1 Error Analysis
% 6.2 Somatization Impact
% 6.3 Evidence Quality
% 6.4 Failure Modes and Abstention

% TODO: Write analysis


\section{Discussion}
\label{sec:discussion}
% Discussion
% - Clinical implications
% - Limitations
% - Ethical considerations
% - Future work

% TODO: Write discussion


\section{Conclusion}
\label{sec:conclusion}
\input{sections/conclusion}

% References
\bibliographystyle{plainnat}
\bibliography{references}

% Appendix
\appendix
\section{Supplementary Materials}
\label{sec:supplementary}
% Supplementary Materials
% A.1 Full ICD-10 Disorder Criteria
% A.2 Somatization Ontology Sample
% A.3 Prompt Templates
% A.4 Detailed Per-Disorder Results
% A.5 Case Studies

% TODO: Write supplementary


\end{document}
